% Generated draft: requires mathematical and source review.
\documentclass[11pt]{article}
\usepackage[margin=25mm]{geometry}
\usepackage{fontspec}
\setmainfont{Latin Modern Roman}
\usepackage{amsmath,amssymb,mathtools,unicode-math}
\setmathfont{Latin Modern Math}
\usepackage{enumitem,needspace,xurl,hyperref}
\setlength{\parindent}{0pt}
\setlength{\parskip}{5pt}
\setlength{\emergencystretch}{3em}
\newcommand{\N}{\mathbb{N}}
\newcommand{\Z}{\mathbb{Z}}
\newcommand{\Q}{\mathbb{Q}}
\newcommand{\R}{\mathbb{R}}
\newcommand{\C}{\mathbb{C}}
\newcommand{\F}{\mathbb{F}}
\newcommand{\A}{\mathbb{A}}
\newcommand{\PP}{\mathbb{P}}
\newcommand{\E}{\mathbb{E}}
\newcommand{\eps}{\varepsilon}
\newcommand{\dd}{\,\mathrm{d}}
\begin{document}
\section*{UnsolvedMath ID: 8800062}
\textbf{Source problem code:} catalogue entry
\subsection{Definitions and mathematical statement}
In the usual Euclidean formulation, let \(L=B\Z^n\subset\R^n\) be a full rank lattice given by a nonsingular binary-encoded integer basis. Its successive minima are
\[
 \lambda_i(L)=\inf\{r>0:\dim_\R\operatorname{span}(L\cap B_2(0,r))\geq i\}.
\]
For a factor \(\gamma(n)\geq1\), approximate SVP asks for \(v\in L\setminus\{0\}\) with \(\|v\|_2\leq\gamma(n)\lambda_1(L)\). Approximate SIVP asks for \(n\) linearly independent vectors \(v_1,\ldots,v_n\in L\) with \(\max_i\|v_i\|_2\leq\gamma(n)\lambda_n(L)\). SIVP does not ask that these vectors form a lattice basis.

Micciancio asks to prove worst-case hardness of both general-lattice tasks for small polynomial approximation factors. A representative form of such a factor is \(\gamma(n)=n^c\) with a small fixed \(c>0\), but the source gives no exponent or exact factor. It does not define hardness: NP-hardness under a specified reduction and conditional exclusion of polynomial-time algorithms are different possible targets. General lattices are not restricted to cyclic, ideal, or NTRU lattices.

The supplied summary recasts the question as small-constant hardness; this does not match the source's small-polynomial-factor target. Its mentioned factor \(2^{(\log n)^{1-\varepsilon}}\), for fixed \(0<\varepsilon<1\), is subpolynomial, not a fixed positive power \(n^c\). Its claim that larger-factor hardness is known but small-constant hardness is not is also logically problematic for the same problem and reduction model: an algorithm achieving a smaller approximation factor also achieves any larger one. No purported resolution or exact hardness threshold from that summary is adopted.
\subsection{Short English statement}
Prove worst-case hardness of approximating general SVP and SIVP within factors that are small powers of the dimension. The source does not specify the exponent, norm, or formal hardness model, and the summary changes the target to constants.
\subsection{Sources}
\begin{itemize}
\item Local UnsolvedMath registry, permanent ID 8800062, source code catalogue entry; attributed to Daniele Micciancio; no proposal year supplied; update 2026-08-06.
\item IPAM, \emph{Workshop on Number Theory and Cryptography --- Open Problems} (2006), attachment 6240, PDF pages 6--7: \url{https://www.ipam.ucla.edu/programs/workshops/workshop-i-number-theory-and-cryptography-open-problems/?tab=schedule}. Registry access date 2026-07-29; PDF and legacy-PPT visual audit. Provenance reports 19 attachments, 17 PDFs and two legacy PPT decks, with repetitions and solved definitions excluded. Present status and redistribution terms require review.
\item The summary names Ajtai (1996), Ajtai--Kumar--Sivakumar, Micciancio, and Khot, without precise papers or reduction hypotheses. Its historical hardness assertions and constant-factor reformulation require correction and verification rather than serving as a resolution of this source problem.
\end{itemize}
\end{document}
